\chapter{Веб-архитектура}
\label{ch:08}

В главе~\ref{ch:07} мы поставили над методом контроллера аннотацию
\texttt{@GetMapping}, запустили приложение и увидели в браузере JSON. Всё
работало. Но если спросить, что именно произошло между нажатием клавиши Enter в
адресной строке и появлением данных на экране, честный ответ большинства
читателей будет таким: «Spring что-то сделал». Эта глава разбирает это
«что-то». Spring~--- не магия, а аккуратная надстройка над двумя вещами:
протоколом HTTP и спецификацией сервлетов. Водить автомобиль, зная только руль
и педали, можно~--- до первой поломки на трассе. Прочитав главу, вы
разберётесь в клиент-серверной архитектуре, научитесь читать и составлять
HTTP-сообщения руками, поймёте, как устроен контейнер сервлетов и жизненный
цикл сервлета, и освоите архитектурный стиль REST~--- фундамент, на котором
стоят все веб-фреймворки.

\section{Клиент-серверная архитектура}

\term{Клиент-серверная архитектура} (\eng{client-server architecture})~---
модель взаимодействия, в которой одна программа (\term{клиент}) отправляет
запросы, а другая (\term{сервер}) принимает их, обрабатывает и отправляет
ответы. Клиент всегда начинает разговор, сервер всегда ждёт.

Представьте ресторан. Вы~--- клиент: садитесь за стол, зовёте официанта,
делаете заказ. Кухня~--- сервер: она не выбегает в зал предлагать котлету, она
ждёт заказа. Вы не знаете, сколько там поваров и откуда привозят мясо; вам
важен результат на тарелке и меню, по которому можно заказывать. Меню здесь~---
интерфейс (в веб-мире это API), тарелка~--- ответ.

Отсюда свойства модели: роли асимметричны, клиент активен, а сервер сам данные
не присылает; один сервер обслуживает много клиентов и потому обязан работать
параллельно~--- здесь пригодится всё, что говорилось о многопоточности в
главе~\ref{ch:05}; обе стороны заранее договорились о протоколе и формате
данных. При этом «сервер»~--- это роль в диалоге, а не железный ящик в
дата-центре: ваш ноутбук со Spring Boot одновременно и сервер, и клиент для
браузера, который к нему стучится.

\subsection{От двух звеньев к трём}

\term{Двухзвенная архитектура} (\eng{two-tier})~--- самый простой вариант:
клиент напрямую общается с сервером базы данных. Так строили корпоративные
системы в 1990-е годы: программа на компьютере бухгалтера сама открывала
соединение с базой и слала туда \texttt{SELECT}. Такого клиента называют
\term{толстым} (\eng{fat client})~--- вся логика живёт в нём. Проблемы вылезают
быстро: при изменении ставки налога программу надо обновить на трёхстах
компьютерах; у клиента есть логин и пароль от базы, поэтому любой сотрудник с
отладчиком получает прямой доступ ко всем таблицам; наконец, триста клиентов
держат триста соединений, и база ложится.

\term{Трёхзвенная архитектура} (\eng{three-tier}) вставляет между клиентом и
базой третье звено~--- \term{сервер приложений}, в котором живёт вся
бизнес-логика. Обе схемы сведены на рис.~\ref{fig:08-1}; сравните, с кем
разговаривает клиент в каждой из них.

\begin{figure}[htbp]\centering
\begin{tikzpicture}[font=\small]
\node[boxw=26mm] (a1){Клиент\\[1pt]\scriptsize толстый};
\node[boxw=30mm,right=20mm of a1] (b1){Сервер БД};
\draw[flowarrow] (a1)--(b1)node[midway,above,font=\scriptsize]{SQL};
\node[boxw=26mm,below=13mm of a1] (a2){Клиент\\[1pt]\scriptsize браузер};
\node[boxw=34mm,right=11mm of a2] (b2)
  {Сервер приложений\\[1pt]\scriptsize Tomcat, Spring Boot};
\node[boxw=26mm,right=11mm of b2] (c2)
  {Сервер БД\\[1pt]\scriptsize PostgreSQL};
\draw[flowarrow] (a2)--(b2)node[midway,above,font=\scriptsize]{HTTP};
\draw[flowarrow] (b2)--(c2)node[midway,above,font=\scriptsize]{JDBC};
\node[anchor=east,font=\scriptsize\itshape]at([xshift=-4mm]a1.west){два звена};
\node[anchor=east,font=\scriptsize\itshape]at([xshift=-4mm]a2.west){три звена};
\end{tikzpicture}
\caption{Двухзвенная и трёхзвенная архитектуры}\label{fig:08-1}
\end{figure}

Вернёмся к ресторану. Двухзвенная архитектура~--- это когда посетитель сам
ходит на склад за продуктами. Трёхзвенная~--- настоящий ресторан: посетитель
общается только с залом, зал передаёт заказ на кухню, кухня берёт продукты со
склада, а на склад посетителя не пускают. Названия трёх уровней и их роли
собраны в табл.~\ref{tab:08-1}; правый столбец показывает, чем каждый уровень
был реализован в предыдущих главах.

\begin{table}[htbp]
\caption{Три звена трёхзвенной архитектуры}
\label{tab:08-1}
\centering
\begin{tabular}{L{3.4cm}L{5.8cm}L{5.6cm}}
\toprule
Звено & Что делает & Чем реализовано в курсе \\
\midrule
Представление (\eng{presentation}) & показывает данные, принимает ввод
  пользователя & браузер, HTML и Thymeleaf, мобильное приложение,
  \texttt{curl} \\
Логика (\eng{application}) & проверяет права, применяет бизнес-правила &
  \texttt{@RestController}, \texttt{@Service} \\
Данные (\eng{data}) & хранит и выдаёт данные & \texttt{@Repository}, JDBC,
  Hibernate, СУБД \\
\bottomrule
\end{tabular}
\end{table}

Выигрыш очевиден: точка обновления одна, пароль от базы знает только сервер
приложений, одно и то же API обслуживает браузер, мобильное приложение и чужую
систему, а при росте нагрузки поднимают несколько копий сервера приложений за
балансировщиком. Клиента, который не содержит логики и умеет только показывать
пришедшее с сервера, называют \term{тонким} (\eng{thin client}): его не нужно
обновлять на каждой машине, зато он не работает без сети. Современные
одностраничные приложения~--- гибрид: код исполняется в браузере, но
доставляется с сервера, а за данными ходит в REST API. Поэтому большая часть
этой главы~--- про данные и REST: сервер всё чаще отдаёт не готовый HTML, а
JSON.

\section{Веб-сервер, сервер приложений и контейнер сервлетов}

\term{Веб-сервер} (\eng{web server})~--- программа, которая принимает
HTTP-запросы и отдаёт содержимое, обычно статические файлы: HTML, CSS,
JavaScript, картинки. Классические представители~--- Apache HTTP Server и
Nginx. Веб-сервер~--- это газетный киоск: у продавца на полках лежат готовые
газеты, вы называете номер, он достаёт нужную и отдаёт; он ничего не пишет и не
готовит. Кроме раздачи файлов он умеет шифровать соединение, сжимать ответы,
кешировать и распределять нагрузку между приложениями, работая \term{обратным
прокси} (\eng{reverse proxy}). Чего он не умеет~--- выполнять ваш Java-код и
ходить в базу данных.

\term{Сервер приложений} (\eng{application server})~--- платформа, которая
запускает серверные приложения и даёт им среду выполнения: обработку сетевых
соединений, разбор HTTP, управление потоками, пулы соединений, транзакции,
безопасность. Если веб-сервер~--- киоск с готовой прессой, то сервер
приложений~--- кухня ресторана: он готовит ответ под конкретный заказ~---
сходит в базу, посчитает, соберёт JSON и отдаст именно вам. В мире Java
различают два уровня: \term{контейнер сервлетов} (\eng{servlet container})
реализует спецификацию Jakarta Servlet~--- сервлеты, фильтры, слушатели,
сессии; это Apache Tomcat, Eclipse Jetty, Undertow. \term{Полный сервер
приложений} реализует всю платформу Jakarta EE, добавляя JPA, CDI, JTA, JMS и
EJB; это WildFly, Payara, WebSphere, WebLogic. Нам нужен именно контейнер
сервлетов, и он делает за вас всю грязную работу:

\begin{enumerate}
\item открывает серверный сокет и слушает порт (по умолчанию 8080);
\item принимает TCP-соединение и читает из него текст HTTP-запроса;
\item разбирает текст: выделяет метод, путь, заголовки, тело;
\item заворачивает запрос в объект \texttt{HttpServletRequest}, а заготовку
      ответа~--- в \texttt{HttpServletResponse};
\item находит по адресу нужный сервлет и вызывает его метод в потоке из пула;
\item превращает записанный ответ обратно в текст HTTP и шлёт в сокет;
\item закрывает или переиспользует соединение.
\end{enumerate}

Именно поэтому вы ни разу не писали \texttt{new ServerSocket(8080)} в
Spring-проекте: за вас это сделал встроенный Tomcat. В промышленной
эксплуатации веб-сервер и сервер приложений работают в связке: снаружи стоит
Nginx на порту~443, он терминирует HTTPS и сам отдаёт статические файлы, а всё,
что нужно посчитать, передаёт Tomcat на порт~8080, который, в свою очередь,
ходит в базу данных по JDBC.

Исторически Java-приложение упаковывали в \term{WAR-архив} (\eng{Web
Application Archive}) со строго определённой структурой
(рис.~\ref{fig:08-2}) и клали в каталог \texttt{webapps} уже установленного
Tomcat.

\begin{figure}[htbp]\centering
\begin{forest} hierarchy
[myapp.war
  [WEB-INF/
    [web.xml -- дескриптор развёртывания]
    [classes/ -- скомпилированные .class]
    [lib/ -- jar-зависимости]
  ]
  [index.html]
  [css/style.css]
]
\end{forest}
\caption{Структура WAR-архива}\label{fig:08-2}
\end{figure}

Spring Boot перевернул схему: теперь Tomcat встраивается внутрь приложения, и
на выходе получается обычный исполняемый JAR, который запускается командой
\texttt{java -jar myapp.jar}. Разница как между «привезти мебель в готовую
квартиру» и «привезти дом-вагончик вместе с мебелью»: второе проще перевозить и
разворачивать в облаке~--- поэтому сегодня побеждает именно этот подход.

\section{Протокол HTTP}

\term{HTTP} (\eng{HyperText Transfer Protocol})~--- прикладной протокол
передачи данных по модели «запрос~--- ответ». Работает поверх TCP (в
HTTP/3~--- поверх QUIC и UDP), по умолчанию занимает порт~80, а защищённая
версия HTTPS~--- порт~443. Всё остальное определяют три его свойства. Протокол
\emph{текстовый} (до HTTP/2): запрос и ответ~--- буквально строки, которые
можно набрать руками. Протокол работает \emph{без состояния}: сервер ничего не
помнит между запросами. И протокол \emph{расширяемый}: всё, чего не хватает,
добавляется заголовками~--- так появились кеширование, авторизация и сжатие.

Аналогия~--- обычная почта: каждый запрос это письмо в конверте, на котором
написано, кому и что вы хотите (стартовая строка), в углу стоят служебные
пометки «срочно» и «формат A4» (заголовки), а внутри лежит лист с содержанием
(тело). Почтальон не помнит, о чём было предыдущее письмо: каждое письмо
самодостаточно.

\subsection{Структура запроса и ответа}

Запрос состоит ровно из четырёх частей, идущих строго в таком порядке:
стартовая строка вида «метод, путь, версия»; заголовки~--- строки вида «имя:
значение», по одной на строку; пустая строка-разделитель, обязательная;
тело~--- необязательное, оно есть у \texttt{POST}, \texttt{PUT} и
\texttt{PATCH} и обычно отсутствует у \texttt{GET}. Разберите по этим четырём
частям запрос из листинга~\ref{lst:08-1}: найдите, где кончаются заголовки и
начинается тело.

\begin{listing}[H]
\begin{minted}{text}
POST /api/books HTTP/1.1
Host: library.example.com
User-Agent: curl/8.4.0
Accept: application/json
Content-Type: application/json; charset=UTF-8
Content-Length: 70

{"title":"Война и мир","author":"Толстой","year":1869}
\end{minted}
\caption{HTTP-запрос целиком}
\label{lst:08-1}
\end{listing}

Заголовок \texttt{Host} в HTTP/1.1 обязателен: на одном IP-адресе живут сотни
сайтов, и сервер по нему выбирает нужный. Обратите внимание и на
\texttt{Content-Length}: в теле 54~символа, а в заголовке стоит 70~--- каждая
кириллическая буква в UTF-8 занимает два байта, и длина считается в байтах, а
не в символах. Посчитаете символами~--- сервер оборвёт тело на середине или
будет ждать недостающие байты до тайм-аута. Ещё деталь: строки в HTTP
разделяются последовательностью CR~LF, а не одним переводом строки.

Структура ответа зеркальная и отличается только первой строкой: вместо
стартовой идёт \term{строка состояния} вида «версия, код, пояснение». В
листинге~\ref{lst:08-2} сервер сообщает, что создал ресурс, и указывает его
адрес заголовком \texttt{Location}~--- по нему клиент узнаёт, какой
идентификатор присвоила база.

\begin{listing}[H]
\begin{minted}{text}
HTTP/1.1 201 Created
Location: /api/books/42
Content-Type: application/json; charset=UTF-8
Content-Length: 78

{"id":42,"title":"Война и мир","author":"Толстой","year":1869}
\end{minted}
\caption{HTTP-ответ на запрос создания книги}
\label{lst:08-2}
\end{listing}

\subsection{Из чего состоит адрес}

Адрес ресурса раскладывается на шесть частей. Разберём строку
\texttt{https://site.example.com:8443/books/42?lang=ru\#top}~--- результат
разбора приведён в табл.~\ref{tab:08-2}, и особое внимание обратите на
последнюю строку.

\begin{table}[htbp]
\caption{Части адреса ресурса}
\label{tab:08-2}
\centering
\begin{tabular}{L{3.4cm}L{5.2cm}L{6.2cm}}
\toprule
Часть & Значение в примере & Кто её видит \\
\midrule
Схема & \texttt{https} & определяет протокол и порт по умолчанию \\
Хост & \texttt{site.example.com} & уходит на сервер в заголовке
  \texttt{Host} \\
Порт & \texttt{8443} & если не указан~--- 80 для \texttt{http}, 443 для
  \texttt{https} \\
Путь & \texttt{/books/42} & уходит в стартовую строку запроса \\
Строка запроса & \texttt{lang=ru} & тоже в стартовую строку, после знака
  вопроса \\
Фрагмент & \texttt{\#top} & на сервер не отправляется вообще, обрабатывается
  браузером \\
\bottomrule
\end{tabular}
\end{table}

Отсюда практическое правило: никогда не передавайте пароли и токены в строке
запроса~--- они попадут в журналы веб-сервера, в историю браузера и в заголовок
\texttt{Referer} при переходе на чужой сайт. Секретам место только в теле
запроса или в заголовке \texttt{Authorization}.

Строка запроса и тело HTML-формы кодируются по одним правилам; формат
называется \texttt{application/x-www-form-urlencoded}: пробел записывается
знаком плюса, а символы со специальным смыслом~--- амперсанд, знак равенства,
решётка~--- процентной последовательностью \texttt{\%26}, \texttt{\%3D} и
подобными, иначе сервер принял бы амперсанд внутри значения за разделитель
параметров. По той же причине в стартовую строку нельзя класть кириллицу: там
разрешены только символы ASCII, и слово «Иван» уходит на сервер как
\texttt{\%D0\%98\%D0\%B2\%D0\%B0\%D0\%BD}~--- по два байта на букву.

\subsection{HTTP своими руками}

Лучший способ поверить, что HTTP~--- это просто текст, написать сервер на голом
сокете и посмотреть, что прилетает от браузера. Всё нужное для этого разобрано
в главе~\ref{ch:05}: серверный сокет, потоки ввода-вывода, кодировки. Сервер из
листинга~\ref{lst:08-3} печатает в консоль стартовую строку и заголовки
запроса, а ответ склеивает вручную из четырёх частей.

\begin{listing}[H]
\begin{minted}{java}
package ru.fa.web;

// import java.io.*; java.net.*; java.nio.charset.StandardCharsets;

public class RawHttpServer {

    public static void main(String[] args) throws IOException {
        try (ServerSocket server = new ServerSocket(8080)) {
            System.out.println("Слушаю порт 8080, стоп по Ctrl+C");
            while (true) {
                // accept() ждёт, пока кто-нибудь не подключится
                try (Socket socket = server.accept()) {
                    handle(socket);
                }
            }
        }
    }

    private static void handle(Socket socket) throws IOException {
        BufferedReader in = new BufferedReader(new InputStreamReader(
                socket.getInputStream(), StandardCharsets.UTF_8));
        // Заголовки читаем до первой пустой строки
        String line;
        while ((line = in.readLine()) != null && !line.isEmpty()) {
            System.out.println(line);
        }
        // Ответ собираем сами: код, заголовки, пустая строка, тело
        byte[] body = "<h1>Привет из сокета</h1>"
                .getBytes(StandardCharsets.UTF_8);
        String head = "HTTP/1.1 200 OK\r\n"
                + "Content-Type: text/html; charset=UTF-8\r\n"
                + "Content-Length: " + body.length + "\r\n"
                + "Connection: close\r\n\r\n";
        OutputStream out = socket.getOutputStream();
        out.write(head.getBytes(StandardCharsets.US_ASCII));
        out.write(body);
        out.flush();
    }
}
\end{minted}
\caption{Простейший HTTP-сервер на голом сокете}
\label{lst:08-3}
\end{listing}

Запустите класс и откройте в браузере адрес \texttt{http://localhost:8080/}~---
в консоли появится настоящий запрос браузера со всеми его заголовками. Здесь
нет ни одной строки бизнес-логики, только разбор текста и склейка ответа: ровно
эту работу за вас делает Tomcat, но на пару порядков аккуратнее. И заметьте,
что этот сервер обслуживает клиентов строго по одному: второй посетитель ждёт,
пока уйдёт первый. Именно за это мы и платим контейнеру сервлетов.

\section{Методы HTTP}

\term{Метод} (\eng{HTTP method}), он же глагол, говорит серверу, что мы хотим
сделать с ресурсом. Ресурс задан путём (существительное), метод~--- действие
(глагол): запись \texttt{GET /books/42} читается как «прочитать книгу
номер~42». Семь основных методов и два их важнейших свойства сведены в
табл.~\ref{tab:08-3}.

\begin{table}[htbp]
\caption{Методы HTTP и их свойства}
\label{tab:08-3}
\centering
\begin{tabular}{L{2.2cm}L{5.6cm}C{1.8cm}C{2.2cm}C{2.2cm}}
\toprule
Метод & Назначение & Тело запро\-са & Безопас\-ный & Идемпо\-тентный \\
\midrule
\texttt{GET} & получить представление ресурса & нет & да & да \\
\texttt{POST} & создать подчинённый ресурс или запустить обработку & да & нет &
  нет \\
\texttt{PUT} & полностью заменить ресурс по известному адресу & да & нет & да \\
\texttt{PATCH} & частично изменить ресурс & да & нет & нет \\
\texttt{DELETE} & удалить ресурс & нет & нет & да \\
\texttt{HEAD} & то же, что \texttt{GET}, но только заголовки & нет & да & да \\
\texttt{OPTIONS} & узнать, какие методы поддерживает ресурс & нет & да & да \\
\bottomrule
\end{tabular}
\end{table}

Бытовая аналогия~--- библиотечная картотека. \texttt{GET}~--- посмотреть
карточку книги, не трогая её. \texttt{POST}~--- принести в библиотеку новую
книгу, инвентарный номер ей присвоит библиотекарь. \texttt{PUT}~--- заменить
карточку номер~42 целиком новым бланком, \texttt{PATCH}~--- зачеркнуть в ней
одну строку, \texttt{DELETE}~--- списать книгу из фонда. \texttt{HEAD}~---
спросить, есть ли карточка и когда её правили, не читая содержимого: так
проверяют, не изменился ли большой файл, прежде чем качать его целиком.

Разница между \texttt{POST} и \texttt{PUT}~--- самая частая ошибка на экзамене.
\texttt{POST} отправляется на \emph{коллекцию} (\texttt{POST /api/books}), и
идентификатор назначает сервер. \texttt{PUT} отправляется на \emph{конкретный
ресурс} (\texttt{PUT /api/books/42}) и означает «пусть по этому адресу лежит
вот это»; если ресурса не было, \texttt{PUT} может его создать, но адрес всё
равно выбрал клиент.

\term{Безопасный метод} (\eng{safe}) не изменяет состояние ресурса на сервере,
это только чтение: такие запросы можно повторять и кешировать.
\term{Идемпотентный метод} (\eng{idempotent})~--- тот, многократное повторение
которого оставляет сервер в том же состоянии, что и однократное; речь идёт
именно о состоянии сервера, а не о коде ответа. Идемпотентность~--- это кнопка
вызова лифта: нажмёте один раз или десять, лифт приедет один раз. А кнопка
«купить билет» не идемпотентна: десять нажатий~--- десять билетов и десять
списаний с карты, поэтому такие кнопки и блокируют после первого щелчка.

Почему \texttt{DELETE} идемпотентен, хотя изменяет данные? После первого
\texttt{DELETE /books/42} книги нет; после второго её тоже нет~--- состояние
сервера одно и то же, а разные коды ответа (204 и 404) состояния не меняют.
Почему \texttt{PATCH} неидемпотентен? Потому что патч бывает разный: тело
«установить статус: выдана» повторять безопасно, а тело «увеличить счётчик
выдач на единицу» меняет данные при каждом повторе. Практический вывод: браузер
и прокси могут автоматически повторить \texttt{GET} или \texttt{PUT} при обрыве
связи, а повторять \texttt{POST} они не станут~--- иначе пользователь дважды
оформит заказ.

\section{Коды состояний}

\term{Код состояния} (\eng{status code})~--- трёхзначное число, первая цифра
которого задаёт класс ответа: \texttt{1xx}~--- «принял, продолжаю»,
\texttt{2xx}~--- «готово», \texttt{3xx}~--- «иди туда», \texttt{4xx}~--- «ты
неправ», \texttt{5xx}~--- «я сломался». Граница между двумя последними классами
важнее всех остальных. Код \texttt{4xx} означает: запрос неисправим в текущем
виде, повторять его бессмысленно, чините клиент. Код \texttt{5xx} означает:
запрос был нормальный, это сервер не справился~--- можно повторить позже. Если
приложение отдаёт 500 в ответ на неверный ввод пользователя, оно врёт клиенту.
Коды, которые нужно знать наизусть, собраны в табл.~\ref{tab:08-4}.

\begin{table}[htbp]
\caption{Коды состояний, которые встречаются каждый день}
\label{tab:08-4}
\centering
\begin{tabular}{L{1.2cm}L{4.4cm}L{9.4cm}}
\toprule
Код & Название & Когда возвращать \\
\midrule
200 & OK & успешный \texttt{GET}, \texttt{PUT}, \texttt{PATCH} \\
201 & Created & успешный \texttt{POST}, создавший ресурс; обязателен заголовок
  \texttt{Location} \\
204 & No Content & успех, но тела нет; классика для \texttt{DELETE} \\
304 & Not Modified & ресурс не менялся, берите из кеша; тела нет \\
400 & Bad Request & сервер не понял запрос: битый JSON, нет параметра \\
401 & Unauthorized & вы не представились или токен неверен \\
403 & Forbidden & мы знаем, кто вы, но вам сюда нельзя \\
404 & Not Found & ресурса по такому адресу нет \\
405 & Method Not Allowed & ресурс есть, но метод не поддерживается \\
415 & Unsupported Media Type & сервер не умеет разбирать присланный
  \texttt{Content-Type} \\
422 & Unprocessable Content & синтаксис верный, но данные не проходят
  бизнес-проверку \\
500 & Internal Server Error & необработанное исключение в коде сервера \\
\bottomrule
\end{tabular}
\end{table}

Три пары кодов путают чаще всего. \emph{401 против 403}: первый означает
«покажите документы», второй~--- «документы в порядке, но вход только для
сотрудников». \emph{400 против 422}: 400~--- сервер не смог разобрать
сообщение, 422~--- разобрал прекрасно, но содержимое нарушает правила
предметной области, например год издания 3025. \emph{404 против 405}: 404~---
нет такого адреса, 405~--- адрес есть, но глагол не тот. Запомните главное: код
ответа~--- часть вашего API, и если клиент проверяет только, пришёл ли JSON, он
будет считать ошибку успехом. Поэтому в REST-контроллере главы~\ref{ch:07} мы
возвращали \texttt{ResponseEntity}, а не голый объект.

\section{Заголовки, состояние и сессии}

Заголовки~--- это метаданные сообщения. Их сотни, но рабочий минимум невелик и
приведён в табл.~\ref{tab:08-5}; средний столбец подсказывает, в запросе
заголовок встречается или в ответе.

\begin{table}[htbp]
\caption{Рабочий минимум заголовков HTTP}
\label{tab:08-5}
\centering
\begin{tabular}{L{4.2cm}L{1.8cm}L{9.0cm}}
\toprule
Заголовок & Где & Назначение \\
\midrule
\texttt{Host} & запрос & какой из виртуальных хостов нужен; обязателен \\
\texttt{Accept} & запрос & какой формат ответа устроит клиента \\
\texttt{Content-Type} & оба & в каком формате тело этого сообщения \\
\texttt{Content-Length} & оба & размер тела в байтах, а не в символах \\
\texttt{Authorization} & запрос & учётные данные: токен или строка
  \texttt{Basic} \\
\texttt{Cookie}, \texttt{Set-Cookie} & оба & вернуть и выдать cookie \\
\texttt{Location} & ответ & адрес созданного (201) или нового (3xx) ресурса \\
\texttt{Allow} & ответ & разрешённые методы; обязателен при 405 \\
\bottomrule
\end{tabular}
\end{table}

Классическая ошибка новичка~--- перепутать \texttt{Content-Type} и
\texttt{Accept}. \texttt{Content-Type} описывает то, что вы отправляете прямо
сейчас; \texttt{Accept}~--- то, что хотите получить в ответ. В GET-запросе
\texttt{Content-Type} не нужен вовсе: тела-то нет. Механизм, при котором клиент
просит формат, а сервер выбирает подходящий, называется \term{согласованием
содержимого} (\eng{content negotiation}); если сервер не умеет ничего из
запрошенного, он отвечает кодом 406.

\subsection{Почему HTTP вас не помнит}

\term{Отсутствие состояния} (\eng{statelessness})~--- фундаментальное свойство
HTTP: сервер не хранит контекст между запросами, и каждый запрос обязан
содержать всю информацию, необходимую для его обработки. Представьте официанта
с полной амнезией: вы говорите «мне суп»~--- он приносит; вы говорите «и то же
самое другу»~--- а он не помнит ни вас, ни супа. Жить с таким официантом можно
двумя способами: либо каждый раз повторять весь заказ целиком (это REST с
токеном в каждом запросе), либо получить номерок и предъявлять его при каждом
обращении (это cookie и сессии). Зачем нужно такое неудобное свойство? Оно даёт
масштабируемость: если сервер ничего не помнит, любой из десяти его клонов
обработает любой запрос. Как только сервер начинает помнить пользователя в
собственной памяти, вы обязаны возвращать его на ту же машину, а при её падении
теряете все данные.

\term{Cookie}~--- небольшой фрагмент данных, который сервер просит клиента
сохранить и присылать обратно при каждом следующем запросе к этому сайту.
Cookie~--- это номерок из гардероба: сам он ничего не значит, ценность в том,
что по нему гардеробщик находит ваше пальто. Из атрибутов важнее прочих три:
\texttt{HttpOnly} делает cookie недоступной из JavaScript и защищает от кражи
через XSS, \texttt{Secure} разрешает отправку только по HTTPS, а
\texttt{SameSite} ограничивает отправку с чужих сайтов и защищает от CSRF.
\term{Сессия} (\eng{session})~--- область данных на сервере, связанная с
пользователем и найденная по идентификатору из cookie. При входе сервер создаёт
сессию, кладёт в неё имя и роль пользователя и присылает заголовок
\texttt{Set-Cookie} со значением \texttt{JSESSIONID=8F3A21}; при следующем
запросе браузер возвращает этот идентификатор, и сервер узнаёт посетителя. В
сервлетах работа с сессией укладывается в пять строк (листинг~\ref{lst:08-4}).

\begin{listing}[H]
\begin{minted}{java}
HttpSession session = request.getSession();  // создаст, если нет
session.setAttribute("user", "ivan");        // положить
String user = (String) session.getAttribute("user");
session.setMaxInactiveInterval(30 * 60);     // жизнь 30 минут
session.invalidate();                        // выход: уничтожить
\end{minted}
\caption{Работа с сессией в сервлете}
\label{lst:08-4}
\end{listing}

Сессии нарушают идею \eng{stateless}, но не на уровне протокола, а на уровне
приложения: протокол по-прежнему передаёт всё нужное в каждом запросе, а сервер
уже помнит состояние. Именно поэтому REST их не любит: при десяти экземплярах
сервера сессия, созданная на первом, не видна остальным. Решения два~--- общее
хранилище сессий (например, Redis) либо самодостаточные токены JWT из
главы~\ref{ch:07}.

Осталось сказать про защиту. \term{HTTPS}~--- тот же HTTP, но внутри
шифрованного канала \term{TLS} (\eng{Transport Layer Security}): сначала
устанавливается TCP-соединение, затем проводится рукопожатие, в котором сервер
предъявляет сертификат и стороны договариваются о ключах, и только потом внутрь
туннеля идут обычные HTTP-запросы. TLS даёт конфиденциальность, целостность и
аутентификацию сервера. HTTP/2 сохранил всю семантику, но заменил текст
двоичными кадрами и добавил мультиплексирование, а HTTP/3 заменил TCP на
протокол QUIC поверх UDP. Для нас важно одно: методы, коды, заголовки и REST от
версии протокола не зависят~--- меняется упаковка, смысл остаётся.

\section{Сервлеты}

\term{Сервлет} (\eng{servlet})~--- Java-класс, который работает внутри
контейнера сервлетов и обрабатывает запросы клиентов по протоколу HTTP. Это
стандарт, описанный спецификацией Jakarta Servlet (раньше~--- Java Servlet,
пакет \texttt{javax.servlet}; начиная с Jakarta~EE~9~---
\texttt{jakarta.servlet}). Сервлет~--- это арендатор офиса в бизнес-центре: вы
не строите здание, не проводите электричество и не нанимаете охрану~--- всё это
даёт вам центр, то есть контейнер, а вы въезжаете в комнату с готовой розеткой
и занимаетесь своим делом, соблюдая правила центра.

Правила описаны интерфейсом \texttt{jakarta.servlet.Servlet}, но напрямую его
почти никто не реализует: обычно наследуются от класса \texttt{HttpServlet},
который уже умеет разбирать HTTP и раскладывать запросы по методам
\texttt{doGet}, \texttt{doPost}, \texttt{doPut} и \texttt{doDelete}.
Унаследованные реализации этих методов не пустые: если вы метод не
переопределили, сработает вариант из \texttt{HttpServlet}, отвечающий кодом 405.
Само API подключают зависимостью \texttt{jakarta.servlet-api} с областью
видимости \texttt{provided}: реализацию даёт контейнер.

\subsection{Жизненный цикл сервлета}

Контейнер управляет сервлетом по строгому сценарию из трёх этапов
(рис.~\ref{fig:08-3}). Обратите внимание на петлю в середине: обслуживание
повторяется на каждый запрос, а всё остальное происходит по одному разу за
жизнь приложения.

\begin{figure}[htbp]\centering
\begin{tikzpicture}[font=\small,node distance=7mm]
\node[boxw=24mm] (c){\texttt{new}\\[1pt]\scriptsize объект создан};
\node[boxw=28mm,right=of c] (i){\texttt{init()}\\[1pt]\scriptsize один раз};
\node[boxw=40mm,right=of i] (v)
  {\texttt{service()}\\[1pt]\scriptsize \texttt{doGet}, \texttt{doPost}, ...};
\node[boxw=32mm,below=16mm of v] (d)
  {\texttt{destroy()}\\[1pt]\scriptsize один раз};
\node[term,text width=32mm,left=12mm of d] (g){отдан сборщику мусора};
\draw[flowarrow] (c)--(i);
\draw[flowarrow] (i)--(v);
\draw[flowarrow] (v.north east)..controls +(9mm,11mm) and +(-9mm,11mm)..
  (v.north west)node[midway,above,font=\scriptsize]{на каждый запрос};
\draw[flowarrow] (v)--(d)
  node[midway,right,font=\scriptsize]{остановка приложения};
\draw[flowarrow] (d)--(g);
\end{tikzpicture}
\caption{Жизненный цикл сервлета}\label{fig:08-3}
\end{figure}

В методе \texttt{init()} читают конфигурацию и готовят дорогие ресурсы, в
\texttt{doGet()} и его родственниках обрабатывают запрос, в \texttt{destroy()}
освобождают ресурсы. Отсюда три вывода. Первый: экземпляр сервлета один на всё
приложение~--- создавать новый объект на каждый запрос было бы слишком дорого.
Второй: сервлет обязан быть потокобезопасным, потому что все запросы идут через
один и тот же объект в разных потоках пула. Поле \texttt{private int counter} в
сервлете~--- гарантированная гонка данных из главы~\ref{ch:05}; изменяемое
состояние держите в локальных переменных метода, в сессии или в
потокобезопасных структурах вроде \texttt{AtomicInteger}. Третий: по умолчанию
сервлет создаётся лениво, при первом обращении; чтобы контейнер создал его
сразу при старте, указывают параметр \texttt{loadOnStartup}.

\subsection{Объекты запроса и ответа}

\texttt{HttpServletRequest}~--- объект, в который контейнер сложил разобранный
HTTP-запрос; это ваш единственный источник знаний о том, чего хочет клиент. У
него спрашивают метод (\texttt{getMethod}), путь (\texttt{getRequestURI}),
хвост пути после шаблона сервлета (\texttt{getPathInfo}), параметр
(\texttt{getParameter}), заголовок (\texttt{getHeader}), тело
(\texttt{getReader} или \texttt{getInputStream}) и сессию
(\texttt{getSession}). \texttt{HttpServletResponse}~--- объект, через который вы
собираете ответ: \texttt{setStatus} задаёт код, \texttt{setHeader} и
\texttt{setContentType} управляют заголовками, \texttt{getWriter} и
\texttt{getOutputStream} дают символьный и байтовый потоки тела, а
\texttt{sendRedirect} отправляет перенаправление.

С этими объектами связаны четыре ловушки, дающие самые загадочные ошибки.
Первая: тело запроса читается ровно один раз и ровно одним способом~--- вызвать
сначала \texttt{getReader()}, а потом \texttt{getInputStream()} нельзя, второй
вызов бросит \texttt{IllegalStateException}. Вторая: \texttt{getParameter()} для
формы, присланной методом \texttt{POST}, сам вычитывает тело, и после обращения
к параметрам тела вы уже не получите. Третья: сначала статус и заголовки, потом
тело~--- как только буфер сброшен, заголовки ушли клиенту, и менять их поздно.
Четвёртая: кодировку задавайте до первой записи вызовом
\texttt{setContentType} с указанием \texttt{charset}, иначе получите
нечитаемые символы вместо русского текста.

\subsection{Сервлет целиком}

Соберём всё сказанное в один работающий класс (листинг~\ref{lst:08-5}).
Проследите в нём три шага: сначала читаются данные запроса, затем
устанавливаются код и заголовки ответа, и только потом пишется тело.

\begin{listing}[H]
\begin{minted}{java}
package ru.fa.web;

// import jakarta.servlet.http.*; java.io.*;
// import java.util.concurrent.atomic.AtomicInteger;

@WebServlet(urlPatterns = "/hello", loadOnStartup = 1)
public class HelloServlet extends HttpServlet {

    // Экземпляр сервлета один на всё приложение, поэтому счётчик
    // обязан быть потокобезопасным
    private final AtomicInteger hits = new AtomicInteger();

    @Override
    public void init() {
        // Вызывается один раз: здесь готовят дорогие ресурсы
        System.out.println("HelloServlet инициализирован");
    }

    @Override
    protected void doGet(HttpServletRequest request,
                         HttpServletResponse response)
            throws IOException {
        // 1. Читаем данные запроса
        String name = request.getParameter("name");
        if (name == null || name.isBlank()) {
            name = "гость";
        }
        int number = hits.incrementAndGet();

        // 2. Настраиваем ответ: сначала код и заголовки
        response.setStatus(HttpServletResponse.SC_OK);
        response.setContentType("text/html; charset=UTF-8");
        response.setHeader("X-Hit-Number", String.valueOf(number));

        // 3. И только потом пишем тело
        PrintWriter out = response.getWriter();
        out.println("<h1>Здравствуйте, " + name + "!</h1>");
        out.println("<p>Обращение номер " + number + "</p>");
    }

    @Override
    public void destroy() {
        System.out.println("Всего обращений: " + hits.get());
    }
}
\end{minted}
\caption{Сервлет-приветствие со счётчиком обращений}
\label{lst:08-5}
\end{listing}

Зарегистрировать сервлет можно тремя способами: аннотацией \texttt{@WebServlet}
(в Spring Boot для этого нужна аннотация \texttt{@ServletComponentScan} на
классе приложения), дескриптором \texttt{web.xml} и программно, через бин
\texttt{ServletRegistrationBean}. Контейнер выбирает сервлет по шаблону пути в
строгом порядке приоритета: точное совпадение (\texttt{/hello}), префикс пути
(\texttt{/api/books/*}), причём побеждает более длинный, шаблон по расширению
(\texttt{*.pdf}) и в последнюю очередь сервлет по умолчанию на одиночном слеше.
Вот почему в Spring Boot работает всё: \texttt{DispatcherServlet} зарегистрирован
именно на одиночный слеш.

\subsection{Что из этого делает Spring}

Теперь соберём картину целиком (рис.~\ref{fig:08-4}). Читайте схему сверху
вниз: первые блоки~--- работа контейнера, следующие~--- работа фреймворка, и
лишь один блок написан вами.

\begin{figure}[htbp]\centering
\begin{tikzpicture}[font=\small,node distance=4mm]
\node[boxw=40mm] (q){HTTP-запрос};
\node[boxw=80mm,below=of q] (t)
  {\textbf{Tomcat}: разбирает текст, создаёт объекты запроса и ответа};
\node[boxw=80mm,below=of t] (f)
  {\textbf{Фильтры}: логирование, CORS, Spring Security};
\node[boxw=80mm,below=of f] (ds)
  {\textbf{DispatcherServlet}: сервлет на одиночном слеше};
\node[boxw=80mm,below=of ds] (hm)
  {\textbf{HandlerMapping}: какой метод отвечает за этот адрес};
\node[boxw=80mm,below=of hm] (ha)
  {\textbf{HandlerAdapter}: вызов метода с подстановкой аргументов};
\node[boxw=80mm,below=of ha] (ct)
  {\textbf{Ваш метод контроллера}: бизнес-логика};
\node[boxw=80mm,below=of ct] (mc)
  {\textbf{HttpMessageConverter} или \textbf{ViewResolver}: JSON либо HTML};
\node[boxw=40mm,below=of mc] (rs){HTTP-ответ};
\draw[flowarrow] (q)--(t);\draw[flowarrow] (t)--(f);\draw[flowarrow] (f)--(ds);
\draw[flowarrow] (ds)--(hm);\draw[flowarrow] (hm)--(ha);
\draw[flowarrow] (ha)--(ct);\draw[flowarrow] (ct)--(mc);
\draw[flowarrow] (mc)--(rs);
\end{tikzpicture}
\caption{Путь запроса через Spring MVC}\label{fig:08-4}
\end{figure}

Ключевая мысль всей главы: \texttt{DispatcherServlet}~--- обычный сервлет.
Аннотация \texttt{@GetMapping}~--- не новый протокол, а способ сказать этому
сервлету, какой метод вызвать, когда \texttt{request.getMethod()} вернёт строку
\texttt{GET}, а путь совпадёт с шаблоном. До этого уровня можно спуститься
прямо из контроллера, объявив параметрами метода \texttt{HttpServletRequest} и
\texttt{HttpServletResponse}.

\section{Архитектурный стиль REST}

\term{REST} (\eng{Representational State Transfer}, передача состояния
представления)~--- архитектурный стиль построения распределённых систем,
описанный Роем Филдингом в 2000~году. Филдинг был одним из авторов спецификации
HTTP, и REST~--- по сути обобщение принципов, на которых уже был построен веб.
Ключевое слово~--- \emph{стиль}, а не протокол и не библиотека: это набор
ограничений, которым система либо соответствует, либо нет; соответствующую им
систему называют \eng{RESTful}. REST~--- это правила дорожного движения, а не
марка автомобиля: они говорят не на чём ехать, а как себя вести, чтобы все
понимали друг друга. И второе заблуждение: REST~--- это не «когда JSON»; о
формате данных REST не говорит ничего.

\subsection{Шесть ограничений}

\begin{enumerate}
\item \textbf{Клиент-сервер.} Клиент занимается интерфейсом, сервер~---
      хранением и логикой; они развиваются независимо.
\item \textbf{Отсутствие состояния.} Каждый запрос содержит всю информацию для
      его обработки; аутентификация~--- в каждом запросе, а не в сессии. Даёт
      горизонтальное масштабирование ценой повторяющихся данных.
\item \textbf{Кешируемость.} Ответ помечается как кешируемый или некешируемый
      заголовками \texttt{Cache-Control} и \texttt{ETag}.
\item \textbf{Единообразие интерфейса.} Самое объёмное ограничение из четырёх
      подпунктов: у каждого ресурса свой адрес; клиент работает не с самим
      объектом, а с его представлением; сообщения самоописываемы~--- метод,
      \texttt{Content-Type} и код ответа несут всё нужное для обработки;
      гипермедиа управляет состоянием приложения, то есть ответ содержит ссылки
      на следующие возможные действия.
\item \textbf{Слоистая система.} Клиент не знает, общается он напрямую с
      сервером или через посредников~--- кеш, балансировщик, шлюз.
\item \textbf{Код по требованию~--- необязательное.} Сервер может передать
      клиенту исполняемый код: классический пример~--- JavaScript.
\end{enumerate}

Если система нарушает любое из первых пяти ограничений, она не \eng{RESTful}.
На практике чаще всего нарушают отсутствие состояния (кладут данные в серверную
сессию) и единообразие интерфейса (изобретают собственные глаголы в адресах).

\subsection{Ресурс, адрес и представление}

\term{Ресурс} (\eng{resource})~--- любая именуемая сущность, которую можно
адресовать: книга, список книг, пользователь, отчёт за март. \term{Адрес}
(\eng{URI})~--- то, по чему ресурс находят: каждый адрес указывает ровно на
один ресурс. \term{Представление} (\eng{representation})~--- конкретный вид
ресурса, передаваемый по сети: документ JSON, XML, страница HTML, файл PDF.

Разница между ресурсом и представлением~--- как между человеком и его
фотографией. Человек~--- ресурс; фотография в паспорте, селфи в мессенджере и
словесное описание в протоколе~--- три разных представления одного и того же
человека. По сети вы передаёте не человека, а его представление; получатель,
изменив описание и прислав его обратно, просит изменить самого человека. Отсюда
и расшифровка названия: \texttt{GET} забирает представление текущего состояния
ресурса, \texttt{PUT} отправляет представление желаемого.

Правила именования адресов теории почти не содержат, зато именно за них вас
будут оценивать на проверке кода (табл.~\ref{tab:08-6}).

\begin{table}[htbp]
\caption{Правила именования адресов ресурсов}
\label{tab:08-6}
\centering
\begin{tabular}{L{4.4cm}L{5.8cm}L{4.6cm}}
\toprule
Правило & Плохо & Хорошо \\
\midrule
Существительные, не глаголы & \texttt{/getBook?id=42} &
  \texttt{GET /books/42} \\
Множественное число для коллекций & \texttt{/book/42} & \texttt{/books/42} \\
Иерархия для вложенности & \texttt{/getReviewsOfBook?id=42} &
  \texttt{/books/42/reviews} \\
Строчные буквы и дефис & \texttt{/BookReviews} & \texttt{/book-reviews} \\
Без расширения формата & \texttt{/books/42.json} & заголовок \texttt{Accept} \\
Фильтры в строке запроса & \texttt{/books/genre/drama} &
  \texttt{/books?genre=drama} \\
Версия в пути или заголовке & \texttt{/books?v=1} & \texttt{/api/v1/books} \\
\bottomrule
\end{tabular}
\end{table}

Строка запроса предназначена для того, что \emph{не идентифицирует} ресурс, а
уточняет выборку: фильтрация, сортировка, постраничная выдача. Критерий
простой: если убрать параметр и запрос всё равно вернёт осмысленный ресурс,
параметру место в строке запроса. Адрес \texttt{/books?genre=drama} без
параметра даст все книги~--- значит, всё правильно; \texttt{/books/42} без
числа уже ничего не адресует~--- значит, число часть пути.

\subsection{Уровни зрелости Ричардсона}

Леонард Ричардсон предложил шкалу из четырёх уровней, показывающую, насколько
интерфейс близок к настоящему REST. \emph{Уровень~0}, «болото POX»,~--- один
адрес и один метод, обычно \texttt{POST}, а что делать, написано в теле: HTTP
служит транспортом для конверта, не более. \emph{Уровень~1}~--- появились
отдельные адреса для сущностей, но метод по-прежнему один. \emph{Уровень~2}~---
используются подходящие методы и осмысленные коды ответов
(\texttt{DELETE /books/42} в ответ на 204); здесь находится подавляющее
большинство реальных интерфейсов, включая тот, что вы писали в
главе~\ref{ch:07}. \emph{Уровень~3}~--- гипермедиа: ответ содержит ссылки на
дальнейшие действия, и клиенту не нужно зашивать адреса в код. Принцип третьего
уровня называется \term{HATEOAS} (\eng{Hypermedia As The Engine Of Application
State}): клиент узнаёт о доступных действиях из самого ответа, а не из
документации, как в листинге~\ref{lst:08-6}.

\begin{listing}[H]
\begin{minted}{text}
{
  "id": 42,
  "title": "Война и мир",
  "status": "available",
  "_links": {
    "self":   { "href": "/api/books/42" },
    "borrow": { "href": "/api/books/42/loans", "method": "POST" }
  }
}
\end{minted}
\caption{Ответ третьего уровня зрелости}
\label{lst:08-6}
\end{listing}

Если книга уже выдана, сервер просто не пришлёт ссылку на выдачу~--- и клиенту
не нужно знать бизнес-правило «выдать можно только доступную книгу», оно
остаётся на сервере. Это в точности то, как работает обычный сайт: вы не
запоминаете адреса страниц, вы ходите по ссылкам. Третий уровень встречается
редко: ответы объёмнее, клиенты сложнее, а разработчику клиента проще прочитать
документацию, чем писать навигацию по ссылкам.

\section{CRUD, REST и HTTP}

\term{CRUD}~--- акроним четырёх базовых операций над хранимыми данными:
\eng{Create} (создать), \eng{Read} (прочитать), \eng{Update} (изменить),
\eng{Delete} (удалить). Это не архитектурный стиль, а минимальный набор
операций, без которого не обходится ни одна информационная система: любая
картотека умеет ровно это~--- завести карточку, посмотреть её, поправить,
выбросить. Прелесть CRUD в том, что он одинаково ложится на все уровни системы:
на SQL, на слой доступа к данным из главы~\ref{ch:06}, на HTTP и на REST
(табл.~\ref{tab:08-7}).

\begin{table}[htbp]
\caption{Соответствие CRUD, SQL и HTTP}
\label{tab:08-7}
\centering
\begin{tabular}{L{2.4cm}L{2.2cm}L{2.2cm}L{3.6cm}L{3.4cm}}
\toprule
CRUD & SQL & Метод & Адрес & Успешный код \\
\midrule
Create & \texttt{INSERT} & \texttt{POST} & \texttt{/api/books} & 201 и
  заголовок \texttt{Location} \\
Read (список) & \texttt{SELECT} & \texttt{GET} & \texttt{/api/books} & 200 \\
Read (один) & \texttt{SELECT} & \texttt{GET} & \texttt{/api/books/42} & 200
  или 404 \\
Update (целиком) & \texttt{UPDATE} & \texttt{PUT} & \texttt{/api/books/42} &
  200 или 204 \\
Update (частично) & \texttt{UPDATE} & \texttt{PATCH} & \texttt{/api/books/42} &
  200 или 204 \\
Delete & \texttt{DELETE} & \texttt{DELETE} & \texttt{/api/books/42} & 204 \\
\bottomrule
\end{tabular}
\end{table}

На уровне Java-кода тем же операциям соответствуют аннотации
\texttt{@PostMapping}, \texttt{@GetMapping}, \texttt{@PutMapping},
\texttt{@PatchMapping} и \texttt{@DeleteMapping}, а в слое доступа к данным~---
методы \texttt{save}, \texttt{findById}, \texttt{findAll} и
\texttt{deleteById}: все слои системы говорят об одном и том же. Обратите
внимание на третью строку таблицы вместе с последней: после успешного
\texttt{DELETE} тот же самый \texttt{GET} вернёт уже 404~--- ресурс перестал
существовать, и сервер честно об этом сообщает, а не отдаёт 200 с пустым телом.

Не всё в предметной области сводится к четырём операциям. Что делать с
действиями «выдать книгу читателю» или «отменить заказ»? Правильный ход~--- не
изобретать глаголы в адресе, а найти в задаче \emph{ресурс}. Выдача книги~---
это сущность: у неё есть читатель, дата выдачи и срок возврата. Значит, выдать
книгу~--- это создать ресурс «выдача» (\texttt{POST /api/books/42/loans}), а
вернуть книгу~--- \texttt{DELETE /api/loans/17}. Такое переосмысление действия
как ресурса~--- главный приём проектирования REST API: поймали себя на глаголе
в адресе~--- спросите, какое существительное вы на самом деле меняете.

\section{REST, SOAP и RPC}

REST~--- не единственный способ построить взаимодействие двух программ.
\term{RPC} (\eng{Remote Procedure Call})~--- стиль, в котором клиент вызывает
функцию на сервере так, будто она локальная, а адрес указывает не на ресурс, а
на процедуру: \texttt{POST /api/getBookById}. Он хорош, когда операции плохо
ложатся на CRUD или когда важна скорость: двоичные варианты вроде gRPC заметно
быстрее JSON поверх HTTP. Минусы~--- единообразия нет, кеширование не работает,
промежуточные узлы не понимают смысла запроса. \term{SOAP} (\eng{Simple Object
Access Protocol})~--- протокол обмена сообщениями в формате XML поверх HTTP:
каждое сообщение~--- конверт с заголовком и телом, все запросы идут методом
\texttt{POST}, а интерфейс формально описан документом WSDL. REST~--- это
короткое сообщение: быстро и без церемоний, но и без гарантий; SOAP~---
заказное письмо с описью вложения и заверенной подписью: медленно, дорого, зато
формально безупречно, поэтому он до сих пор жив в банках и государственных
системах. Вывод: для публичного интерфейса и для браузера берите REST, для
внутреннего общения микросервисов~--- gRPC, а SOAP выбирают вынужденно, когда
так требует контрагент.

\praktikum{}

Понадобятся JDK~17 или новее, утилита \texttt{curl} и текстовый редактор.
Создайте рабочий каталог \texttt{http-lab} и проверьте инструменты командами
\texttt{curl --version} и \texttt{java -version}. Задания выполняются по
порядку: каждое следующее опирается на файлы предыдущего.

Два замечания для Windows, без которых половина команд сломается. В PowerShell
слово \texttt{curl}~--- псевдоним командлета \texttt{Invoke-WebRequest} с
совершенно другими ключами, поэтому везде пишите \texttt{curl.exe}, а команду
набирайте одной строкой: обратная косая черта для переноса там не работает.
Команда \texttt{chcp 65001} переключает консоль на UTF-8, иначе русский текст в
ответах превратится в нечитаемые символы; тело запроса при этом всегда кладите
в файл и передавайте ключом \texttt{-d} со значением \texttt{"@body.json"}.

\subsection*{Задание 1. Сырой обмен и разбор адреса}

Ключ \texttt{-v} показывает и запрос, и ответ целиком: строки со знаком
«больше»~--- то, что отправили вы, строки со знаком «меньше»~--- то, что
прислал сервер. Выполните команду \texttt{curl -v
"https://postman-echo.com/\allowbreak{}get?course=java"} (кавычки вокруг адреса
обязательны) и разметьте вывод: подпишите стартовую строку запроса, заголовки
запроса, пустую строку, строку состояния ответа, пять заголовков ответа и тело
ответа.

Затем возьмите адрес \texttt{https://lib.ru:8443/books/42?g=drama\#top} и
заполните таблицу из шести строк: часть адреса, её значение, уходит ли она на
сервер. Последнюю строку проверьте экспериментом: выполните \texttt{curl -v
"https://postman-echo.com/\allowbreak{}get?token=abc\#secret"} и посмотрите на
стартовую строку запроса.

Ответьте письменно: сколько заголовков \texttt{curl} добавил в запрос
самостоятельно и зачем нужен каждый; попал ли фрагмент в стартовую строку и кто
его отбросил; в каких трёх местах осядет в открытом виде значение
\texttt{token=abc}, даже если соединение защищено, и куда его следовало
положить.

\subsection*{Задание 2. Журнал запросов к эхо-серверу}

Сервис \texttt{postman-echo.com}~--- эхо-сервер: он принимает любой запрос и
возвращает JSON с описанием того, что получил. Это зеркало в примерочной: вы не
гадаете, как на вас сидит рубашка, а смотрите и видите. Создайте файл
\texttt{body.json} с единственной строкой
\texttt{\{"title":"Peace","year":1869\}} в кодировке UTF-8 без BOM и выполните
команды из листинга~\ref{lst:08-7}, записывая для каждой её саму, код ответа и
одну-две строки о том, что нового вы увидели.

\begin{listing}[H]
\begin{minted}{bash}
# параметры в строке запроса и свои заголовки
curl -i "https://postman-echo.com/get?group=PI24&topic=http"
curl -i -H "X-Student: Ivanov" https://postman-echo.com/headers
# тело в формате JSON, затем то же тело как обычный текст
curl -i -X POST "https://postman-echo.com/post" \
     -H "Content-Type: application/json" -d "@body.json"
curl -i -X POST "https://postman-echo.com/post" \
     -H "Content-Type: text/plain" -d "@body.json"
# то же самое в виде формы
curl -i -X POST "https://postman-echo.com/post" \
     --data-urlencode "title=War and Peace" -d "year=1869"
# другие методы, только заголовки, коды состояний
curl -i -X PUT "https://postman-echo.com/put" -d "a=1"
curl -i -X DELETE "https://postman-echo.com/delete"
curl -I "https://postman-echo.com/get"
curl -i "https://postman-echo.com/status/404"
# аутентификация: логин postman, пароль password
curl -i "https://postman-echo.com/basic-auth"
curl -v -u postman:password "https://postman-echo.com/basic-auth"
# cookie, перенаправление и тайм-аут
curl -c cookies.txt "https://postman-echo.com/cookies/set?t=1"
curl -b cookies.txt "https://postman-echo.com/cookies"
curl -i --max-time 1 "https://postman-echo.com/delay/3"
\end{minted}
\caption{Серия запросов к эхо-серверу}
\label{lst:08-7}
\end{listing}

Соберите результаты в сводную таблицу не менее чем из двенадцати строк: номер,
метод, адрес, заголовки, которые задали вы, код ответа, ключевое поле ответа.

Ответьте письменно: в какие три разных поля ответа попали данные при отправке
их как JSON, как формы и как строки запроса; байты в первых двух командах
\texttt{POST} одинаковы~--- почему результат разный и как это объясняет
происхождение кода~415; что вернул сервер на запрос без учётных данных;
является ли Base64 шифрованием и допустима ли такая аутентификация поверх
обычного HTTP.

\subsection*{Задание 3. Свой сервер: от голого сокета до REST}

Внутри \texttt{http-lab} создайте каталоги \texttt{ru/fa/web}~--- они
соответствуют пакету \texttt{ru.fa.web}. Наберите в файл
\texttt{RawHttpServer.java} класс из листинга~\ref{lst:08-3}, соберите его
командой \texttt{javac -encoding UTF-8 -d out ru/fa/web/RawHttpServer.java} и
запустите командой \texttt{java -cp out ru.fa.web.RawHttpServer}. В другом окне
терминала обратитесь к нему через \texttt{curl}, затем откройте тот же адрес в
браузере и сравните два набора заголовков.

Теперь напишите настоящий REST-сервис. Разбирать HTTP руками больше не нужно: в
JDK есть лёгкий сервер \texttt{com.sun.net.httpserver.HttpServer}, не требующий
ни Maven, ни Tomcat. Наберите в файл \texttt{LibraryServer.java} код из
листинга~\ref{lst:08-8}; четыре метода в нём намеренно оставлены
незаполненными.

\begin{listing}[H]
\begin{minted}{java}
package ru.fa.web;

// import com.sun.net.httpserver.*; java.io.*; java.net.*;
// import java.nio.charset.*; java.util.*; java.util.regex.*;
// import java.util.concurrent.*; java.util.concurrent.atomic.*;

public class LibraryServer {

    private static final String BASE = "/api/books";
    // Запросы идут в разных потоках: структура потокобезопасная
    private static final Map<Long, String> STORAGE =
            new ConcurrentHashMap<>();
    private static final AtomicLong ID = new AtomicLong();
    // Наивный разбор JSON. В рабочем коде берут Jackson или Gson
    private static final Pattern TITLE =
            Pattern.compile("\"title\"\\s*:\\s*\"([^\"]*)\"");

    public static void main(String[] args) throws IOException {
        HttpServer server =
                HttpServer.create(new InetSocketAddress(8081), 0);
        server.createContext(BASE, LibraryServer::handle);
        // Пул потоков - то же, что делает контейнер сервлетов
        server.setExecutor(Executors.newFixedThreadPool(8));
        server.start();
    }

    // Разбор пути и метода: работа DispatcherServlet вручную
    private static void handle(HttpExchange ex) throws IOException {
        String method = ex.getRequestMethod();
        // Хвост после /api/books: пустой или вида /42
        String tail = ex.getRequestURI().getPath()
                .substring(BASE.length());
        if (tail.isEmpty() || tail.equals("/")) {
            switch (method) {
                case "GET" -> list(ex);
                case "POST" -> create(ex);
                default -> notAllowed(ex, "GET, POST");
            }
            return;
        }
        long id;
        try {
            id = Long.parseLong(tail.substring(1));
        } catch (NumberFormatException e) {
            send(ex, 400, "{\"error\":\"нужно число\"}");
            return;
        }
        switch (method) {
            case "GET" -> readOne(ex, id);
            case "DELETE" -> delete(ex, id);
            default -> notAllowed(ex, "GET, DELETE");
        }
    }

    private static void create(HttpExchange ex) throws IOException {
        String type = ex.getRequestHeaders().getFirst("Content-Type");
        if (type == null || !type.startsWith("application/json")) {
            send(ex, 415, "{\"error\":\"нужен application/json\"}");
            return;
        }
        String body = new String(ex.getRequestBody().readAllBytes(),
                StandardCharsets.UTF_8);
        Matcher m = TITLE.matcher(body);
        if (!m.find()) {
            // Разобрано, но не проходит проверку данных
            send(ex, 422, "{\"error\":\"нужно поле title\"}");
            return;
        }
        long id = ID.incrementAndGet();
        STORAGE.put(id, m.group(1));
        // 201 обязан сопровождаться заголовком Location
        ex.getResponseHeaders().set("Location", BASE + "/" + id);
        send(ex, 201, "{\"id\":" + id + "}");
    }

    // Отправка ответа: сначала код и заголовки, потом тело
    private static void send(HttpExchange ex, int code, String data)
            throws IOException {
        byte[] body = data.getBytes(StandardCharsets.UTF_8);
        ex.getResponseHeaders().set("Content-Type",
                "application/json; charset=UTF-8");
        ex.sendResponseHeaders(code, body.length);
        try (OutputStream out = ex.getResponseBody()) {
            out.write(body);
        }
    }
}
\end{minted}
\caption{Каркас REST-сервиса «Библиотека»}
\label{lst:08-8}
\end{listing}

Допишите четыре метода: \texttt{list} отдаёт кодом 200 массив всех книг;
\texttt{readOne}~--- одну книгу кодом 200 или, если её нет, кодом 404;
\texttt{delete} удаляет книгу и отвечает кодом 204 без тела, для чего длину тела
передают равной $-1$; \texttt{notAllowed} ставит заголовок \texttt{Allow} со
списком разрешённых методов и отвечает кодом 405.

Запустите сервер и во втором терминале выполните полный цикл жизни ресурса:
получите пустую коллекцию, создайте книгу (ожидается 201 и заголовок
\texttt{Location}), прочитайте её (200), удалите (204), прочитайте снова (404).
После этого намеренно сломайте четыре вещи и запишите код каждого ответа:
отправьте \texttt{POST} без заголовка \texttt{Content-Type}; отправьте тело без
поля \texttt{title}; выполните \texttt{DELETE} на коллекцию
\texttt{/api/books}; запросите \texttt{/api/books/abc}. Оформите результаты
таблицей «проверка, метод, адрес, ожидаемый и полученный коды».

Ответьте письменно: какая строка кода отвечает за одновременное обслуживание
нескольких клиентов и чем это отличается от \texttt{RawHttpServer}; почему
хранилище объявлено как \texttt{ConcurrentHashMap}; какие коды вернулись на два
запроса \texttt{DELETE} подряд и идемпотентен ли \texttt{DELETE}; почему на
запрос без заголовка \texttt{Content-Type} сервер отвечает 415, а на тело без
поля \texttt{title}~--- 422.

\subsection*{Задание 4. Сервлет в контейнере сервлетов}

Теперь запустим сервлет внутри настоящего контейнера и посмотрим, что берёт на
себя Tomcat: серверный сокет, разбор текста запроса, пул потоков, жизненный
цикл объекта. Создайте проект Spring Boot с единственной зависимостью
\eng{Spring Web}~--- она приносит встроенный Tomcat и API сервлетов. Рядом с
аннотацией \texttt{@SpringBootApplication} поставьте
\texttt{@ServletComponentScan}: без неё Spring Boot не найдёт классы,
помеченные \texttt{@WebServlet}. Перенесите в проект класс
\texttt{HelloServlet} из листинга~\ref{lst:08-5}, изменив только строку
\texttt{package}, запустите приложение и выполните три запроса:
\texttt{curl -i -G "http://localhost:8080/hello" --data-urlencode "name=Иван"},
затем то же без параметра и, наконец,
\texttt{curl -i -X POST "http://localhost:8080/hello"}.

Ответьте письменно: сколько раз после десяти запросов в консоли появилось
сообщение из метода \texttt{init} и что это доказывает про число экземпляров
сервлета; что случилось бы, объяви мы счётчик как \texttt{private int hits};
какой код вернул \texttt{POST} и какой метод класса \texttt{HttpServlet} его
сформировал, если \texttt{doPost} мы не переопределяли.

\subsection*{Задание 5. Проектирование REST API «Библиотека»}

Спроектируйте интерфейс библиотеки с четырьмя сущностями: \emph{книга}
(идентификатор, название, автор, год, жанр, статус), \emph{читатель}
(идентификатор, имя, номер билета), \emph{выдача} (идентификатор, книга,
читатель, дата выдачи, срок возврата), \emph{отзыв} (идентификатор, книга,
читатель, оценка, текст).

Первая часть~--- таблица «операция, метод, адрес, успешный код, возможные коды
ошибок» не менее чем из шестнадцати строк, покрывающая все четыре сущности.
Обязательно включите в неё чтение списка и одной книги, создание, полную замену,
частичное изменение и удаление книги, поиск по жанру, постраничную выдачу,
сортировку по году, отзывы на книгу, регистрацию читателя, а также запрос о
поддерживаемых методах и проверку, изменилась ли книга, без скачивания её
содержимого.

Вторая часть~--- разбор чужих ошибок. Для каждой строки укажите нарушенное
правило и корректный вариант: \texttt{POST /api/getBookById?id=42};
\texttt{GET /api/BookReviews}; \texttt{POST /api/books/42/doDelete};
\texttt{GET /api/book/42.json}; \texttt{GET /api/books/genre/drama};
\texttt{POST /api/books/42/setStatus}.

Третья часть~--- действия, которые не сводятся к четырём операциям. Найдите в
каждом действии существительное и оформите таблицей «действие, найденный
ресурс, метод, адрес, код успеха»: выдать книгу читателю; вернуть книгу;
продлить срок выдачи; получить отчёт за март.

Четвёртая часть~--- проверка на ограничения REST. Для каждого случая укажите
нарушенное ограничение и способ исправления. Первый: после входа сервер кладёт
роль пользователя в сессию и достаёт её оттуда при каждом следующем запросе.
Второй: все ответы отдаются с заголовком \texttt{Cache-Control: no-store},
включая справочник жанров, который меняется раз в год. Третий: чтобы получить
книгу, клиент отправляет \texttt{POST /api/gateway} с телом, где действие
указано полем \texttt{action}. Четвёртый: клиент зашит на конкретный порт
приложения и перестаёт работать при появлении балансировщика.

Ответьте письменно: по какому критерию вы решали, что уходит в путь, а что в
строку запроса; какое ограничение нарушено в каждом из четырёх случаев; на
каком уровне зрелости Ричардсона находится спроектированный вами интерфейс и
что нужно добавить, чтобы поднять его на следующий.

\subsection*{Результаты занятия}

По итогам занятия у вас должны быть готовы: размеченный вывод команды
\texttt{curl -v} и таблица разбора адреса; журнал запросов к эхо-серверу;
файлы \texttt{RawHttpServer.java} и \texttt{LibraryServer.java} с дописанными
методами и таблицей проверки кодов 200, 201, 204, 400, 404, 405, 415 и 422;
проект Spring Boot с работающим сервлетом; четыре таблицы проектирования
интерфейса библиотеки; письменные ответы на вопросы всех пяти заданий.

\kontrolvoprosy

\begin{enumerate}
\item Кто в клиент-серверной архитектуре начинает общение? Какие три проблемы
      двухзвенной архитектуры решает третье звено?
\item Чем веб-сервер отличается от сервера приложений? Что такое контейнер
      сервлетов и какие семь шагов он выполняет между приходом байтов в сокет и
      вызовом вашего кода?
\item Перечислите четыре части HTTP-запроса и четыре части HTTP-ответа. Чем
      первая строка запроса отличается от первой строки ответа?
\item Из каких частей состоит адрес ресурса и какая из них никогда не уходит на
      сервер? Почему в строке запроса нельзя передавать пароль?
\item В чём разница между \texttt{POST} и \texttt{PUT}? Что такое безопасный и
      что такое идемпотентный метод?
\item Почему \texttt{DELETE} считается идемпотентным, хотя изменяет состояние
      сервера, а \texttt{PATCH}~--- не считается?
\item Чем принципиально отличаются классы кодов \texttt{4xx} и \texttt{5xx}? В
      чём разница между 401 и 403, между 400 и 422, между 404 и 405?
\item Чем отличаются заголовки \texttt{Content-Type} и \texttt{Accept}? Почему
      в GET-запросе \texttt{Content-Type} не нужен?
\item Что означает свойство \eng{stateless} и какое преимущество оно даёт? Как
      устроена связка cookie и серверной сессии?
\item Опишите жизненный цикл сервлета. Сколько экземпляров создаёт контейнер и
      какое следствие это имеет для полей класса? Почему тело запроса можно
      прочитать только один раз?
\item Перечислите шесть ограничений REST. Какое из них необязательное и какие
      нарушают чаще всего? Чем ресурс отличается от представления?
\item Сформулируйте пять правил именования адресов. Что попадает в путь, а что
      в строку запроса?
\item Назовите четыре уровня зрелости Ричардсона и приведите таблицу
      соответствия CRUD, SQL, методов HTTP и кодов ответов.
\end{enumerate}

\testzadaniya

\begin{enumerate}

\item Из каких частей и в каком порядке состоит HTTP-запрос?
  \begin{enumerate}
    \item[а)] стартовая строка, заголовки, пустая строка, тело
    \item[б)] заголовки, стартовая строка, пустая строка, тело
    \item[в)] стартовая строка, тело, пустая строка, заголовки
    \item[г)] строка состояния, заголовки, тело, пустая строка
  \end{enumerate}

\item В адресе \texttt{https://site.com:8443/books/42?lang=ru\#top} какая часть
      не отправляется на сервер?
  \begin{enumerate}
    \item[а)] порт \texttt{8443}
    \item[б)] путь \texttt{/books/42}
    \item[в)] строка запроса \texttt{lang=ru}
    \item[г)] фрагмент \texttt{\#top}
  \end{enumerate}

\item Чем \texttt{POST} принципиально отличается от \texttt{PUT}?
  \begin{enumerate}
    \item[а)] \texttt{POST} создаёт ресурс, а \texttt{PUT} только читает его
    \item[б)] \texttt{POST} передаёт данные в строке запроса, \texttt{PUT}~---
          в теле
    \item[в)] \texttt{POST} отправляется на коллекцию и идентификатор назначает
          сервер, \texttt{PUT}~--- на ресурс, адрес которого выбрал клиент
    \item[г)] \texttt{POST} идемпотентен, а \texttt{PUT}~--- нет
  \end{enumerate}

\item Почему \texttt{DELETE} считается идемпотентным, хотя изменяет данные?
  \begin{enumerate}
    \item[а)] после первого удаления ресурса нет и после второго тоже нет:
          состояние сервера одинаково, хотя код ответа может отличаться
    \item[б)] он не удаляет запись, а лишь помечает её как скрытую
    \item[в)] контейнер сервлетов кеширует ответы на \texttt{DELETE}
    \item[г)] спецификация запрещает отправлять \texttt{DELETE} повторно
  \end{enumerate}

\item Какой код и какой заголовок обязан вернуть сервер после успешного
      создания ресурса методом \texttt{POST}?
  \begin{enumerate}
    \item[а)] 200 OK и заголовок \texttt{Content-Location}
    \item[б)] 204 No Content без дополнительных заголовков
    \item[в)] 202 Accepted и заголовок \texttt{Retry-After}
    \item[г)] 201 Created и \texttt{Location} с адресом нового ресурса
  \end{enumerate}

\item Присланный JSON разобран без ошибок, но год издания книги указан как
      3025. Какой код точнее описывает ситуацию?
  \begin{enumerate}
    \item[а)] 400 Bad Request
    \item[б)] 422 Unprocessable Content
    \item[в)] 409 Conflict
    \item[г)] 415 Unsupported Media Type
  \end{enumerate}

\item Что означает свойство \eng{stateless} у протокола HTTP?
  \begin{enumerate}
    \item[а)] сервер не хранит данные на диске
    \item[б)] сервер не хранит контекст клиента между запросами, поэтому любая
          копия сервера обработает любой запрос
    \item[в)] клиент не сохраняет ничего, включая cookie
    \item[г)] соединение закрывается после каждого запроса
  \end{enumerate}

\item В сервлете объявлено поле \texttt{private int counter}, а в
      \texttt{doGet} выполняется \texttt{counter++}. Что произойдёт под
      нагрузкой?
  \begin{enumerate}
    \item[а)] ничего: контейнер создаёт свой экземпляр на каждый запрос
    \item[б)] поле будет сбрасываться в ноль после каждого запроса
    \item[в)] значения будут теряться: экземпляр один на всё приложение, и
          запросы идут через него в разных потоках~--- это гонка данных
    \item[г)] контейнер выбросит исключение: такие поля запрещены
  \end{enumerate}

\item В методе \texttt{doPost} сначала вызвали \texttt{request.getReader()} и
      прочитали тело, а затем \texttt{request.getInputStream()}. Что
      произойдёт?
  \begin{enumerate}
    \item[а)] код отработает: методы читают из независимых буферов
    \item[б)] второй вызов выбросит \texttt{IllegalStateException}: тело читают
          только одним способом и только один раз
    \item[в)] второй вызов вернёт \texttt{null}, а ответ окажется пустым
    \item[г)] код не скомпилируется
  \end{enumerate}

\item Какое место занимает \texttt{DispatcherServlet} в обработке запроса?
  \begin{enumerate}
    \item[а)] это обычный сервлет на одиночном слеше; он находит нужный метод
          контроллера и вызывает его
    \item[б)] это отдельный сетевой сервер, работающий вместо Tomcat
    \item[в)] это фильтр, выполняемый до всех сервлетов
    \item[г)] это аннотация, которой помечают контроллеры
  \end{enumerate}

\item Какой адрес лучше всего соответствует правилам именования в REST?
  \begin{enumerate}
    \item[а)] \texttt{GET /getReviewsOfBook?id=42}
    \item[б)] \texttt{POST /books/42/doDeleteReview}
    \item[в)] \texttt{GET /BookReviews/42}
    \item[г)] \texttt{GET /books/42/reviews}
  \end{enumerate}

\end{enumerate}
